\chapter{ОГЛЯД ЛІТЕРАТУРИ}

%Hidden cite for correct bibliography ordering
\nocite{bahvalov-et-al,tumor-thermal-model,lung-tumor-thermal-conductivity,database-tissue-properties} 

% https://nauka.gov.ua/information/osnovni-etapy-v-napysanni-ohliadu-literatury-dlia-dysertatsiinoi-roboty/

\section{Аналіз дослідження}

Нормативи оформлення:

1) Кожен абзац літературного огляду подається у формі порівняльного аналізу або резюме того, чи іншого висновку із
першоджерела з припущеннями чи висновками інших авторів (особливо важливо виявити протиріччя у висновках різних авторів
на вирішення однієї і тієї ж проблеми; "посваривши" таким чином двох чи більше авторів і, на цій підставі, зробити свій
висновок; доцільно підтвердити існування невирішеної проблеми та актуальність її вирішення).

2)Кожен абзац літературного огляду, якщо він не відображає власну думку, а надає інформацію з інших першоджерел, повинен
завершуватися посиланням на конкретне першоджерело, яке повинно бути у розділі «Список використаних джерел»

Повнота цієї частини дослідження досягається шляхом пошуку і завантаження тематичного матеріалу, який у свою чергу вже
має власний список літератури та огляд. Кожна позиція повинна бути вивчена і проаналізована.

Пошук нових джерел припиняється, якщо після опрацювання та аналізу безлічі статей в наступних матеріалах не
зустрічаються нові публікації та прізвища. Наступним етапом стає побудова змістовної частини, розподіл за авторськими
концепціями та школами.

\section{Ситуація з науковими школами}

\section{Значущі результати використаних джерел}

\section{Висновок}

